% !TEX root = ../main.tex
\section*{附录 A\quad 支撑材料文件列表}

支撑材料压缩包按以下结构组织，所列程序均为论文建模、求解与复核所用的完整源文件：

\begin{enumerate}
  \item \texttt{AI工具使用详情.pdf}：AI工具名称、用途、使用方式及人工核验说明；
  \item \texttt{源程序/问题一/quest1.py}：半平面交、旋转卡壳与覆盖判定程序；
  \item \texttt{源程序/问题二/q2\_solver\_robust\_20260911.py}：候选源区域、接收保证域及第二检测点鲁棒优化程序；
  \item \texttt{源程序/问题三/}：多频道自动搜索、四叉树可行域维护、离线仿真程序及参数文件；
  \item \texttt{结果数据/q2\_result.json}：问题二复核结果；\texttt{结果数据/q3\_current\_36\_runs.json}：问题三当前参数的 36 局离线运行记录；
  \item \texttt{README\_支撑材料说明.md}：运行环境、入口命令和文件说明。
\end{enumerate}

\noindent 三问程序均仅依赖 Python 标准库。问题三正式入口优先读取同目录下的 \texttt{params.json}，离线复核入口为 \texttt{auto\_test.py}。所有路径均为压缩包内相对路径。

\newcommand{\codeinput}[2]{%
  \subsection*{#1}
  \VerbatimInput[
    fontsize=\tiny,
    numbers=left,
    numbersep=4pt,
    frame=single,
    breaklines=true,
    breakanywhere=true,
    tabsize=4
  ]{#2}%
}

\clearpage
\section*{附录 B\quad 问题一完整源程序}
\codeinput{quest1.py}{appendix_code/问题一/quest1.py}

\clearpage
\section*{附录 C\quad 问题二完整源程序}
\codeinput{q2\_solver\_robust\_20260911.py}{appendix_code/问题二/q2_solver_robust_20260911.py}

\clearpage
\section*{附录 D\quad 问题三完整源程序}
\codeinput{algorithm.py}{appendix_code/问题三/algorithm.py}
\codeinput{auto\_test.py}{appendix_code/问题三/auto_test.py}
\codeinput{basic\_sim.py}{appendix_code/问题三/basic_sim.py}
\codeinput{iterator.py}{appendix_code/问题三/iterator.py}
\codeinput{main.py}{appendix_code/问题三/main.py}
\codeinput{quest3.py}{appendix_code/问题三/quest3.py}
\codeinput{sss.py}{appendix_code/问题三/sss.py}
\codeinput{params.json}{appendix_code/问题三/params.json}
\codeinput{candidates.json}{appendix_code/问题三/candidates.json}
